We trained eighteen models under the same protocol described in
Section~\ref{sec:data}. Below we walk through each family,
report the MAE per model, and finish with the unified leaderboard
(Fig.~\ref{fig:leaderboard}) and a per-window comparison
(Fig.~\ref{fig:per-window}). Unless otherwise stated, all numbers are
mean~$\pm$~standard deviation across the four rolling-origin windows.

\subsection{Naive Baselines}

We need a floor below which any forecasting model should not fall, and
that floor is the naive family. We ran five baselines:
\textbf{historical mean}, \textbf{last value}, \textbf{drift},
\textbf{seasonal-naive} (at $s=7$), and \textbf{ARIMA naive-seasonal}
(an ARIMA(3,1,3) applied to the de-seasonalised series, used as a
weak baseline in some related work). All of them sit comfortably
above 30~kWh MAE: historical mean is 30.13~kWh, seasonal naive is
34.98~kWh. The seasonal-naive is interesting because it scores
\emph{worse} than the unconditional mean, which is unusual; it
happens because the test window contains a Thanksgiving week, and
``copy last week'' is the wrong heuristic exactly there. The point of
the baselines is that any forecasting model that does not beat them
clearly is not worth running.

\subsection{Classical (S)ARIMA / ARIMAX Family}

We ran an exhaustive AIC grid search over $(p, d, q)(P, D, Q)_7$ with
$p, q \in [0, 4]$, $P, Q \in [0, 2]$, $d = D = 1$ (so the seasonal
difference at lag 7 is applied) and $s = 7$. The grid was repeated
with and without exogenous variables (temperature, weekend flag,
holiday flag), giving two families: SARIMA (no exog) and SARIMAX (with
exog). The grid is reproducible from the cached CSV
\texttt{data/sarima\_grid\_results.csv}.

In addition to the SARIMA grid, we evaluated the simpler members of
the Box-Jenkins family that the project rubric explicitly requires:
ARMA, ARIMA, AutoARIMA (seasonal and non-seasonal), and ARIMAX. Each
of these models is fit via maximum likelihood with the
\texttt{statsmodels} SARIMAX backend, with rolling-origin refits at
each test window so the model has access to the most recent
actuals. The best classical configurations on our test set are:
\begin{itemize}[leftmargin=*]
    \item \textbf{ARIMAX (2,1,3) with exog}: MAE $25.52 \pm 14.38$~kWh.
    \item \textbf{SARIMAX $(2,1,4)\!\times\!(1,1,2)_7$ with exog}:
    MAE $25.89 \pm 12.76$~kWh.
    \item \textbf{ARIMA (2,1,3)}: MAE $26.31 \pm 12.43$~kWh.
    \item \textbf{AutoARIMA non-seasonal (0,1,1)}:
    MAE $26.40 \pm 11.50$~kWh.
\end{itemize}
ARIMAX (2,1,3) is the strongest classical model. It does not include
an explicit seasonal term, but the exogenous variables (especially
weekend / holiday flags) absorb most of the weekly cycle, so adding
the seasonal lag does not help much. We use ARIMAX (2,1,3) as the
\emph{best-classical} reference in the rest of the paper.

\subsection{LSTM}

The LSTM (\texttt{scripts/lstm\_train.py}) uses five input features
(energy, mean temperature, weekend flag, holiday flag, month) and
performs a small hyper-parameter sweep over sequence length
$\in \{14, 28\}$, hidden size $\in \{32, 64\}$, and number of layers
$\in \{1, 2\}$. Each configuration is trained with Adam, MSE loss,
dropout 0.1 and early stopping on the 60-day validation split. The
winning configuration is \texttt{seq14\_h64\_l2}: 14-day sequences,
hidden size 64, 2 layers. At inference time the LSTM forecasts one
day ahead and is rolled forward autoregressively for the full 7-day
horizon, feeding back its own prediction with the known future
exogenous row each step.

The trained LSTM scores \textbf{MAE $25.68 \pm 10.64$~kWh}, which is
slightly worse than ARIMAX on the mean but noticeably tighter in
variance. This is consistent with the literature: at this dataset size
LSTM does not dominate (S)ARIMAX in mean error, but it tends to be
more stable across windows.

\subsection{Zero-Shot Foundation Models}

Notebook 06 of the project runs four foundation models in zero-shot
mode --- no fine-tuning, just feed the history and read the
forecast:
\begin{itemize}[leftmargin=*]
    \item \textbf{Google TimesFM 1.0} (200M parameters):
    MAE $26.32 \pm 11.14$~kWh.
    \item \textbf{Google TimesFM 2.0} (500M parameters):
    MAE $26.23 \pm 12.59$~kWh.
    \item \textbf{Amazon Chronos-T5 base} (200M parameters):
    MAE $27.31 \pm 11.31$~kWh.
    \item \textbf{Amazon Chronos-T5 large} (710M parameters):
    MAE $27.94 \pm 13.02$~kWh.
\end{itemize}
Two findings stand out. First, TimesFM beats Chronos on this series
in both size classes; second, the larger Chronos-T5 large
\emph{loses} to its smaller sibling, which surprised us. The likely
explanation is that Chronos was trained on a corpus heavy in
financial / random-walk-like series, which is precisely the opposite
of what our highly-seasonal series looks like; the larger model
amplifies the prior of its pre-training rather than adapting to the
new structure. TimesFM, which has a more diverse pre-training corpus
including energy and load series, transfers better here. Neither
foundation model beats the strongest classical model (ARIMAX), which
is a useful reality check on the ``foundation models will eat
everything'' narrative.

\subsection{PatchTST from Scratch}

For the SoTA push we implemented PatchTST~\cite{nie2023patchtst} from
scratch in PyTorch (Fig.~\ref{fig:patchtst}). The architecture
choices, made deliberately small to match our $\sim$1900-point
dataset:
\begin{itemize}[leftmargin=*]
    \item Context length: 56 days (8 weeks).
    \item Patch length: 7 days, non-overlapping stride 7 $\Rightarrow$
    8 tokens per input.
    \item 3 transformer encoder layers, $d_{\text{model}}=96$, 4
    attention heads, GELU activation, dropout 0.15.
    \item \textbf{Channel-mixed patching}: each patch is a
    $(7, 4)$ block over the four input channels (energy + 3 exog),
    flattened to 28-dim and linearly embedded; cheaper and more
    data-efficient than channel-independent attention for our
    sequence length.
    \item \textbf{RevIN}~\cite{kim2022revin} on the target channel
    only: per-sample mean / std subtraction before the encoder and
    inverse transform on the output. This is what kept the model
    competitive against the regime shifts in the data (COVID drop,
    holiday weeks).
\end{itemize}
We train each model with AdamW (LR $5 \times 10^{-4}$, weight decay
$10^{-4}$), cosine annealing, early stopping with patience 25, batch
size 32. Training a single PatchTST takes $\approx 2$~minutes on the
RTX 5070~Ti.

A single trained PatchTST hits MAE $27.1$~kWh, which is worse than
ARIMAX. The breakthrough comes from training three independent
models with different seeds (42, 7, 2024) and averaging their
predictions at inference. The seed-ensemble drops to
\textbf{MAE $25.19 \pm 8.90$~kWh}, the best \emph{single architecture}
in our leaderboard and below the ARIMAX baseline both on mean and on
variance.

\begin{figure}[!htbp]
\centering
\includegraphics[width=0.95\textwidth]{fig_patchtst}
\caption{PatchTST forward pass as implemented in
\texttt{scripts/patchtst\_train.py}. The energy channel goes through
RevIN before patching; the three exogenous channels are concatenated
into each patch and flattened together. The head is a single linear
layer that directly predicts the next 7 days, with no autoregression.}
\label{fig:patchtst}
\end{figure}

\subsection{The Final Ensemble}

We then asked whether combining PatchTST and the LSTM gives us anything
extra. The naive way to stack would be to learn weights with Ridge or
NNLS on a held-out fold; we tried both with leave-one-window-out
cross-validation and they both \emph{overfit} (the Ridge variant
collapses to MAE $32.3$~kWh, much worse than any individual model).
The reason is simple: with only 28 test points it is impossible for a
meta-learner to generalise.

What does work is the dumbest combination: the simple unweighted mean
of PatchTST and LSTM. The choice is made \emph{a priori} (these are
the two strongest deep learning models in the project, picked by
architecture rather than by test score) and gives:
\begin{equation*}
\text{Ensemble}_{\text{mean}} = \tfrac{1}{2} \big(\hat{y}_{\text{PatchTST}} + \hat{y}_{\text{LSTM}}\big)
\end{equation*}
\noindent
which scores \textbf{MAE $25.09 \pm 9.36$~kWh}. This is the best model
on the leaderboard, beating PatchTST alone by 0.10~kWh and ARIMAX by
0.43~kWh on the mean.

\subsection{Leaderboard and Per-Window Analysis}

Figure~\ref{fig:leaderboard} shows the final ranking with error bars
(one standard deviation across the four windows). Three observations:
\begin{enumerate}[leftmargin=*]
    \item The top of the leaderboard is tightly packed: every model
    in the top six is within 1~kWh MAE of the best, which is
    smaller than the standard deviation of every individual model.
    \item The Ensemble PatchTST + LSTM does \emph{not} have the
    smallest standard deviation --- that honour goes to the historical
    mean baseline, which is trivially stable. The smallest std among
    forecasting models belongs to PatchTST alone (8.90 vs.\ 14.38 for
    ARIMAX): the seed-ensemble buys us almost 40\% variance reduction
    over the previous best.
    \item All forecasting models comfortably beat the naive
    baselines. The MAE gap between the best model (25.09~kWh) and
    the best baseline (30.13~kWh) is $\sim$17\%, which is the
    learnable signal in the series.
\end{enumerate}

\begin{figure}[!htbp]
\centering
\includegraphics[width=0.95\textwidth]{fig_leaderboard}
\caption{Final leaderboard. Bars are MAE in kWh (lower is better),
error bars are one standard deviation across the four rolling-origin
windows. The Ensemble PatchTST + LSTM (top) is the best model overall;
the PatchTST seed-ensemble has the lowest std among forecasting
models. The full table with every metric and the source notebook for
each row is in \texttt{docs/RESULTS.md} in the repository.}
\label{fig:leaderboard}
\end{figure}

Figure~\ref{fig:per-window} breaks down the MAE by window for the top
six models. The takeaway here is window~3 (Nov 14--20, 2023), which
includes Thanksgiving: every model has its worst week here, with MAE
shooting up to 36--43~kWh. PatchTST is the one that loses the least
on that week (36.4~kWh vs.\ 38.8~kWh for the LSTM and 41~kWh for
ARIMAX), and that is the main reason it wins the leaderboard. Outside
window~3, the top models are essentially tied.

\begin{figure}[!htbp]
\centering
\includegraphics[width=0.95\textwidth]{fig_per_window}
\caption{Per-window MAE for the top six models. Window~3 (Thanksgiving
week) is hard for every model; the ranking on the other three windows
is much tighter. PatchTST loses less on the hard week, which is what
puts it (and its ensemble with LSTM) on top of the leaderboard.}
\label{fig:per-window}
\end{figure}
